\chapter{Implementation}
The previous chapter explained the architectural design and workflow of each phase.
This chapter explains how the software implements those designs.
The first section of this chapter covers the project setup: the dependencies, local LLM communication and model assignment.
The second section details state management and the GUI architecture.
The remaining sections break down non-trivial components of the system.

\section{Setup and Local LLM Integration}

\paragraph{Dependencies}
The system requires Python 3.11 or newer.
A \texttt{requirements.txt} file lists all Python package dependencies for automated installation.
A LM~Studio\footnote{\url{https://lmstudio.ai/}, accessed 2026-03-05} installation is needed to serve the local LLMs and a local \LaTeX{} distribution is required to compile the final PDF.
The \texttt{lmstudio} Python SDK\footnote{\url{https://github.com/lmstudio-ai/lmstudio-python}, accessed 2026-03-05} (version 1.5+) handles all local LLM communication.
This SDK provides access to local LLMs and embedding models running on the LM~Studio server.
\texttt{pydantic}\footnote{\url{https://github.com/pydantic/pydantic}, accessed 2026-03-05} is used for schema validation of the structured JSON outputs generated by the local LLMs.
The GUI uses \texttt{tkinter}~\cite{Tkinter} and \texttt{sv\_ttk}\footnote{\url{https://github.com/rdbende/Sun-Valley-ttk-theme}, accessed 2026-03-05}, a third-party theme for styling.
\texttt{pillow}\footnote{\url{https://github.com/python-pillow/Pillow}, accessed 2026-03-05} loads template PNG icons and converts them into \texttt{tkinter}-compatible image objects for use in interface buttons and header bars.
It resizes these images to match the active font size and recolors their pixels to fit the current light or dark theme.
\texttt{tkinterweb}\footnote{\url{https://github.com/Andereoo/TkinterWeb}, accessed 2026-03-05} renders HTML blocks.
The GUI uses it to display compiled Markdown instead of raw strings.
The \texttt{pymupdf4llm}\footnote{\url{https://github.com/pymupdf/pymupdf4llm}, accessed 2026-03-05} library converts research paper PDFs into structured Markdown.
\texttt{PyMuPDF}\footnote{\url{https://github.com/pymupdf/PyMuPDF}, accessed 2026-03-05} converts generated PDF plots into PNG images for rendering them in the GUI and for VLM analysis.
The \texttt{requests} library manages HTTP calls to the Semantic Scholar~\cite{kinney2025semanticscholaropendata}, Unpaywall~\cite{Unpaywall} and arXiv~\cite{arXiv} APIs.
The system installs a default scientific stack for the experimentation phase.
If needed, the local LLM generating the experiment can import \texttt{numpy}~\cite{Harris_2020} and \texttt{scipy}~\cite{SciPy} to execute numerical and statistical computations, \texttt{pandas}~\cite{Pandas} to manipulate structured datasets and \texttt{matplotlib} and \texttt{seaborn}~\cite{Bisong2019} to generate the plots to visualize the results.

\paragraph{Local LLM Integration}
LM~Studio handles all LLM inference locally and exposes a localhost API.
The \texttt{lmstudio} SDK manages communication with this API to replace raw HTTP requests.
An interaction typically requires three steps.
A Python script builds a prompt string.
It passes this string to the \texttt{model.respond()} method.
The system then receives a response object containing the generated text.
For embeddings, the system calls the \texttt{embedding\_model.embed()} method instead.
This method accepts a list of strings and returns a list of float vectors.

At startup, the system checks for a running inference server by calling a function the LM Studio SDK provides.
If this check fails, the application displays a warning dialog and exits after the user closes it, since each phase requires a local LLM.

If reasoning models are used, they might generate internal thinking blocks.
LM~Studio can read these blocks using \texttt{<think>...</think>} tags.
A post-processing function runs a regular expression to strip these blocks from the response before saving the final output, so the models thinking process does not pollute the result.

\paragraph{Model Configuration}
The \texttt{Settings} class defines 13 variables to assign specific models to certain tasks rather than using one global model.
Table~\ref{tab:model_variables} groups these variables by the phases of the system.

\begin{table}[htbp]
    \centering
    \caption{Model Configuration Variables}
    \label{tab:model_variables}
    \begin{tabularx}{\textwidth}{l l X}
        \toprule
        \textbf{Phase} & \textbf{Variables} & \textbf{Type} \\
        \midrule
        Context Analysis & 2 & Language Model \\
        Paper Search & 1 + 1 & Language Model + Embedding Model \\
        Hypothesis & 1 & Language Model \\
        Experimentation & 4 + 1 & Language Model + Vision-Language Model \\
        Paper Writing & 1 + 1 & Language Model + Embedding Model \\
        \LaTeX{} Generation & 1 & Language Model \\
        \bottomrule
    \end{tabularx}
\end{table}

This design lets the user pick specialized models.
For example, they can assign a model specialized in code generation to the experiment phase and a model specialized in academic text generation to the writing phase.
The experimentation phase requires a VLM to validate generated plots and write their captions.
Embedding models are used to index the parsed papers, rank search results, and retrieve evidence chunks during the writing phase.

Any model variable can be changed during runtime via the settings screen of the GUI.
During saving, the system rewrites the \texttt{settings.py} file directly.
It reads the file, matches the target line with a regular expression, updates the value and writes the file back to disk.
This direct overwrite keeps the settings editable by hand without requiring a separate configuration parser.

\paragraph{Lazy Model Loading}
Loading LLMs consumes VRAM and initialization time.
Loading the model for every variable simultaneously at startup could overwhelm the VRAM of the user's GPU.
Hence, the system loads models on demand via a mixin pattern.

The \texttt{LazyModelMixin} class checks if the requested model exists.
On its first access, it calls \texttt{lms.llm()} with the model name to load it.
Following accesses will only return the loaded instance without reloading the model.
The \texttt{LazyEmbeddingMixin} repeats this logic for embedding models using \texttt{lms.embedding\_model()}.

This on-demand loading works well with LM~Studio's ``Only Keep Last JIT Loaded Model'' option.
When enabled, LM~Studio unloads the previous model before loading the next, so exactly one model is kept in memory.
Executing a new phase triggers a new model request.
If a different model is requested, LM~Studio automatically frees the current one from memory.

The paper writing phase requires an exception because it runs a language model and an embedding model simultaneously.
During the Draft-Critique-Retrieve-Improve pipeline, the system alternates between these two models.
The retrieve step constantly queries the embedding model to search the literature index.
Reloading models on every cycle would cause unnecessary delays.
To solve this problem, the \texttt{LMSJITSettings} context manager writes to LM~Studio's settings file and disables JIT unloading temporarily.
Once writing finishes, it restores the setting and unloads all models from that session.

\paragraph{Directory Architecture}
There are four main directories in the project.
The \texttt{gui/} directory contains all user interface code.
This includes the main application class, screens, popups and dialogs.
The \texttt{phases/} directory contains the backend logic.
It is split into subdirectories for each phase: \texttt{context\_analysis}, \texttt{paper\_search}, \texttt{hypothesis\_\allowbreak{}generation}, \texttt{experimentation}, \texttt{paper\_writing}, and \texttt{document\_compilation}.
These phase modules run independently and never import from the GUI.
The \texttt{utils/} directory holds helper functions for file input and output, LLM calls, model loading and PDF processing.
The \texttt{latex\_templates/} directory stores the \LaTeX{} templates used for document compilation.

All artifacts generated by the system are saved into the \texttt{output/} directory.
To keep the directory organized, the system groups artifacts into three specific subdirectories.
All files related to the Experimentation phase, including Python code, execution logs and plots, are stored in the \texttt{experiments/} folder.
Downloaded PDFs are stored in the \texttt{literature/} folder.
The used \LaTeX{} template and the final PDF are placed in the \texttt{latex/} folder.
Top-level files like \texttt{papers.json} or \texttt{hypothesis.md} store the data for the remaining phases, which only create one or two files each.
User inputs---the paper specification, style guidelines and custom code---are stored in the \texttt{user\_files/} directory.




\section{State Management and User Interface}

This section breaks down the system state and GUI implementation top-down: it starts with the high-level file architecture, continues with the main application window and shared layout blueprint and ends with the individual screens.

\paragraph{File-Based State}
The system writes all generated data to disk as text files.
Text generated by an LLM is saved as Markdown, while structured data is saved as JSON.
The \texttt{file\_utils} module handles all file operations via \texttt{save\_markdown}, \texttt{load\_markdown}, \texttt{save\_json}, and \texttt{load\_json}.

The GUI screens are largely read-only, loading the current state from disk when the user navigates to them.
If the AI produces an error, the user can open the file externally, fix the issue and click ``Reload'' to update the interface.
Pressing the ``Continue'' button changes the view to the next screen, which reads its input directly from the existing files.

\paragraph{Screen Architecture}
The \texttt{PaperGeneratorApp} class inherits from \texttt{tk.Tk} to control the application window.
At startup, it creates a container frame and loads every screen into memory.
It stacks these screens using \texttt{tkinter}'s grid placement.
A \texttt{screen\_order} list specifies the navigation sequence.
The \texttt{show\_frame} method pulls the target screen to the front and runs its \texttt{on\_show} hook.

Every screen of the app inherits from \texttt{BaseFrame}.
The \texttt{BaseFrame} uses a grid layout with five rows:

\begin{enumerate}[nosep]
	\item A header bar with the title and an optional info button.
	\item A horizontal separator line.
	\item A scrollable content area inside a \texttt{tkinter} Canvas.
	\item A bottom separator line.
	\item A navigation bar with navigation and regeneration buttons.
\end{enumerate}

To standardise the UI interactions, \texttt{BaseFrame} provides five hooks:

\begin{enumerate}[nosep]
	\item \textbf{\texttt{create\_content}}: The individual screens override this function to inject text areas, buttons or cards into the Canvas.
	\item \textbf{\texttt{on\_show}}: Runs every time a screen appears. Screens override this to read disk data lazily instead of parsing files during startup. The \texttt{ResearchContextScreen}, for example, reads the context Markdown file only when the user opens the tab.
	\item \textbf{\texttt{on\_next}}: Defines the action for the "Continue" button. Screens override this to trigger conditional logic (like triggering a generation instead of just advancing) before calling the parent method to advance the view.
	\item \textbf{\texttt{on\_back}}: Defines the action for the "Back" button. Screens override this to execute custom logic, like saving settings, before returning to the previous view.
	\item \textbf{\texttt{on\_regenerate}}: Defines the action for the "Regenerate" button. Screens override this to spin up LLM generation in a background thread.
\end{enumerate}

Subclasses configure the \texttt{BaseFrame} during initialization using the following parameters:

\begin{enumerate}[nosep]
	\item \textbf{\texttt{title}}: Sets the header text for the screen.
	\item \textbf{\texttt{has\_next}, \texttt{has\_back}, \texttt{has\_regenerate}}: Boolean flags to toggle the visibility of the respective navigation buttons.
	\item \textbf{\texttt{info\_content}}: Markdown text to display when the user clicks the header's information icon.
	\item \textbf{\texttt{header\_file\_path}}: The path to the artifact file this screen represents. Setting this parameter exposes three action buttons above the main content area.
\end{enumerate}

If a screen defines a \texttt{header\_file\_path}, these action buttons let the user fix data files outside the GUI.
The ``Open in Editor'' button launches the system text editor.
The ``Show in Explorer'' button highlights the file in the OS file manager.
The ``Reload'' button wipes the screen and triggers \texttt{create\_content} again to load the external manual edits.

The \texttt{info\_texts} module stores the documentation for all screens as Markdown strings.
Each screen imports its specific text block and passes it to the \texttt{info\_content} parameter during initialization.
This spawns an information icon in the header bar of the screen.
Clicking the icon opens an \texttt{InfoPopup} displaying the formatted documentation.

\paragraph{Shared UI Components}
The \texttt{BaseFrame.create\_card\_frame} helper handles collapsible UI sections.
It builds a bordered container with a colored header to keep the visual design identical across all screens.
Background threads run slow tasks like paper indexing or LLM requests so the GUI stays responsive.
Triggering the execution of a phase launches a \texttt{ProgressPopup} modal with a spinner and status message.
This popup recursively disables every button in the parent window to prevent duplicate script runs.
The background thread sends status updates to the popup using \texttt{self.after(0, ...)}.
This runs the update safely on the main thread.
Errors trigger a scrollable traceback window inside the modal.

\paragraph{Pipeline Orchestrators}
The GUI uses orchestrator methods to execute an entire phase.
These methods handle both the initial generation and any user-triggered regeneration.
The system maps each of the six phases to a specific core orchestrator:

\begin{itemize}[nosep]
	\item \textbf{Context Analysis}: \texttt{ResearchContextGenerator.generate\_new\_context}
	\item \textbf{Paper Search}: \texttt{LiteratureSearch.run\_automated\_search}
	\item \textbf{Hypothesis Generation}: \texttt{HypothesisBuilder.generate\_new\_hypothesis}
	\item \textbf{Experimentation}: \texttt{ExperimentRunner.run\_new\_experiment}
	\item \textbf{Paper Writing}: \texttt{PaperWritingPipeline.write\_paper}
	\item \textbf{Document Compilation}: \texttt{PaperConverter.generate\_new\_pdf}
\end{itemize}

\paragraph{Theming}
The system allows users to switch between a light and a dark interface to suit their preference and reduce eye strain in low-light environments.
The \texttt{sv\_ttk} Sun Valley theme provides these modes natively.
The app reads the \texttt{Settings.DARK\_MODE} boolean at startup and sets the theme via \texttt{sv\_ttk.set\_theme()}.
Calling \texttt{toggle\_theme} swaps the mode at runtime.
This flips the boolean, re-applies the theme, refreshes all styles, clears the icon cache, and triggers \texttt{apply\_theme\_colors}.

\texttt{sv\_ttk} only restyles \texttt{ttk} widgets.
Standard \texttt{tk} widgets like text areas and border frames require manual repainting.
The \texttt{apply\_theme\_colors} function crawls the widget tree to locate \texttt{tk.Text}, \texttt{tk.Label}, or \texttt{tk.Canvas} instances.
It checks their parent style name and injects the correct hex codes from the \texttt{theme\_colors} module.

\paragraph{Font Sizes}
High-resolution monitors might render default \texttt{tkinter} text unreadably small, or some users might prefer larger text, so the interface allows the user to scale all text elements.
The \texttt{FontManager} class controls the text sizes across headers, buttons, and text fields.
Modifying the font size triggers a broadcast to every registered text element.
This forces \texttt{BaseFrame} to recalculate line wrapping and grid widths.




\section{Literature Search and Selection}

\paragraph{Automated Literature Search}
The \texttt{LiteratureSearch.\allowbreak{}run\_\allowbreak{}automated\_\allowbreak{}search} method is used for the retrieval of academic papers.
First, a local LLM reads the user's research context and generates 15 targeted search queries.
The \texttt{SemanticScholarAPI} class iterates through these queries and uses the Semantic Scholar \texttt{/graph/v1/paper/search} endpoint to retrieve the papers.
The HTTP requests ask for a maximum of 20 results per query alongside specific metadata fields like title, abstract, citation count, open-access status, and DOIs.

The system parses the JSON API responses and converts every valid result into a \texttt{Paper} data object.
This \texttt{dataclass} stores all metadata independently from the API and standardizes the data layout for all later phases.
To clean the raw search data, the system deduplicates the list of papers.
It compares the Semantic Scholar paper IDs and DOIs of the retrieved objects.
As a fallback, it runs a fuzzy string match on the title and first author to drop overlapping entries.

\paragraph{Paper Ranking}
The system ranks the unique \texttt{Paper} objects using a composite score.
The \texttt{PaperRanker.\allowbreak{}rank\_papers} method embeds the title and abstract of each paper.
It computes three metrics:

\begin{enumerate}[nosep]
	\item \textbf{Relevance (80\%):} Cosine similarity aligns the paper's embedding against the research context embedding.
	\item \textbf{Citation Velocity (10\%):} Dividing total citations by age in years produces an age-aware score. Papers with 30+ citations per year score 0.9. Papers older than four years with fewer than 0.5 citations per year receive a strict penalty.
	\item \textbf{Recency (10\%):} An exponential decay function applies a four-year half-life. A paper published today scores 1.0, a four-year-old paper scores 0.5, and an eight-year-old paper scores 0.25.
\end{enumerate}

The final ranking sorts by the weighted sum of these three values.

The \texttt{PaperFilter.filter\_papers} method selects papers in three steps.
First, it auto-selects the top 10 papers by pure relevance, protecting them from deletion.
Second, it fills the remaining positions up to a limit of 40 using the composite score.
Third, it pushes the non-protected papers to an LLM in batches of 10 for filtering via \texttt{PaperFilter.\allowbreak{}verify\_\allowbreak{}with\_\allowbreak{}llm}.
The LLM reads each paper's title, abstract, and the base research context.
It returns a formatted JSON list of the paper numbers to keep.
The prompt instructs the model to delete only papers failing in both method and application.
If the LLM call crashes, the system defaults to keeping the entire batch.

A citation gap detection module scans the final list.
The \texttt{CitationGapFinder.\allowbreak{}identify\_\allowbreak{}missing\_\allowbreak{}papers} method prompts the LLM to identify missing theoretical groundwork.
The \texttt{CitationGapFinder.\allowbreak{}search\_\allowbreak{}suggested\_\allowbreak{}papers} method automatically queries Semantic Scholar to fill these gaps.
It specifically hits the \texttt{/graph/v1/paper/search/match} endpoint to find exact title matches for the LLM suggestions.

\paragraph{Open Access Retrieval and Downloading}
The system needs the physical PDF files to parse their text.
Semantic Scholar results might lack a direct download link.
The \texttt{find\_\allowbreak{}open\_\allowbreak{}access\_\allowbreak{}pdfs} method scans the paper list to find these missing URLs.
It queries the Unpaywall API (\texttt{api.unpaywall.org/v2/\{doi\}}) for any paper with a registered DOI.
It parses the JSON response to extract links to open-access repositories or author websites.
If Unpaywall fails, the system falls back to the arXiv API (\texttt{export.arxiv.org/api/query}).
It runs a title search and parses the XML response.
It applies a Jaccard word similarity threshold of 0.95 between the search query and the arXiv result to prevent incorrect paper matches.
It extracts the PDF link or constructs one from the arXiv ID string.

The \texttt{PDFDownloader} fetches these resolved URLs.
It spoofs a standard browser User-Agent and disables SSL certificate verification to bypass strict server firewalls and 403 blocks.
It saves each downloaded PDF into a subfolder named after the paper ID.




\section{Automated Experimentation}
The experiment runner generates, executes and validates Python code to test the user's hypothesis.

\paragraph{Code Generation and Execution}
The \texttt{ExperimentRunner.\_write\_experiment\_code} method runs a four-step conversational loop.
First, the system initializes an \texttt{lms.Chat} object with the system prompt including the hypothesis, experiment plan, available packages, and the user's custom code files.

The \texttt{CodeExecutor.execute\_file} method runs the generated script in an isolated subprocess.
It calls Python's \texttt{subprocess.Popen} with a configurable 10-minute timeout and captures standard output and standard error.
The executor scans the output directory before and after the run to detect new plot images (PNG, PDF) or data files (JSON, CSV).
It binds these discovered files to the \texttt{ExecutionResult} object.

\paragraph{Error Recovery}
A non-zero return code stops the script and triggers an auto-fix loop in \texttt{ExperimentRunner.\_fix\_experiment\_code}.
The system bundles the broken code, standard output, and standard error into a fix prompt.
It uses a \texttt{Chat} object across attempts so the model reads the history of previous failures.

Timeout errors trigger a special guidance block.
This block forces the model to reduce loop iterations, run in headless mode (e.g., without PyGame UI windows), and simplify matrix operations.
The fix loop runs for up to five attempts before aborting.

\paragraph{Validation and Visual Checking}
A successful execution triggers a separate validation model via \texttt{ExperimentRunner.\_validate\_experiment\_results} to review the code, logs and plots.
The \texttt{ExperimentRunner.\_check\_plots\_visually} method converts any PDF plots to PNG using PyMuPDF.
A VLM scans these PNG images.
It searches for missing axis labels, empty graphs, overlapping text, and missing legends.
The pipeline bundles these visual errors into the primary validation prompt.

The validation model outputs a structured \texttt{ValidationResult} JSON object.
This object flags whether the results are valid and lists specific bugs by line number.
Failed validations bounce back to the code writing model for an automated improvement attempt via \texttt{ExperimentRunner.\_improve\_experiment\_code} before re-executing.

\paragraph{Plot Captioning}
A successful experiment stops the coding loop and the system executes the plot caption generation via \texttt{ExperimentRunner.\_generate\_plot\_captions}.
A VLM receives each plot alongside the hypothesis, experiment plan and standard output log and creates a caption for the image.
The system saves these captions next to the plot files and later injects them into the Results section prompt during the writing phase.




\section{Draft-Critique-Retrieve-Improve Pipeline}
The \texttt{PaperWritingPipeline} class writes the paper sections in a fixed order: Methods, Results, Discussion, Introduction, Related Work, Conclusion, and Abstract.

To build each section, the pipeline loops through four steps based on the PaperQA~\cite{skarlinski2024languageagentsachievesuperhuman} algorithm:

\begin{enumerate}[nosep]
	\item \textbf{Initial Draft:} The \texttt{PaperWriter.generate\_initial\_section} method generates a first version using the paper catalog, research context, experiment logs, and prior sections.
	A system prompt forces the model to insert exact citation keys from the catalog.
	\item \textbf{Critique:} The \texttt{SectionCritic.critique\_section} method reviews the draft.
	It outputs an improvement summary and a JSON list of search queries targeting weak claims.
	The pipeline extracts every citation key from the text and cross-references them against the catalog.
	It flags unrecognized keys to force corrections.
	\item \textbf{Evidence Search:} The system builds a searchable vector index.
	The \texttt{PaperIndexer.\allowbreak{}index\_\allowbreak{}papers} method removes abstracts and conclusions, isolates code blocks, and splits the remaining text into overlapping paragraph chunks of up to 700 tokens.
	A word count heuristic (\texttt{word\_count / 0.75}) estimates block sizes to bypass slow API tokenizer calls.
	The embedding model converts these chunks into vectors and caches them on disk.
	The \texttt{EvidenceGatherer.batch\_search} method embeds the search queries and calculates their cosine similarity against the chunk vectors to retrieve the top matches.
	To refine the selection, a local LLM reads these retrieved chunks in batches.
	It scores their contextual relevance from 0.0 to 1.0 and generates a concise summary for each chunk.
	The prompt forces the model to strip out preexisting in-text citations from the summary.
	Finally, the system deduplicates the evidence and keeps a maximum of two chunks per paper per query.
	\item \textbf{Rewrite:} The \texttt{PaperWriter.rewrite\_section} method ingests the initial draft, the critique feedback, the hallucination warnings, and the new evidence summaries to output the final section text.
\end{enumerate}

Sections requiring no external evidence---the Abstract and Conclusion---bypass the search step.




\section{Citation Management and \LaTeX{} Assembly}

\paragraph{Citation Extraction}
The \texttt{generate\_literature\_bib} function scans the final Markdown draft for citation keys formatted as \texttt{[citationKey]} using regular expressions.
It cross-references each found key against the \texttt{Paper} objects stored in memory.
It extracts the metadata---title, authors, year, DOI, journal---and creates a \texttt{bibliography.bib} file.
If a key fails to match an indexed paper, it inserts a fallback placeholder entry to prevent compilation failures.

\paragraph{\LaTeX{} Template System}
The \texttt{latex\_templates/} directory can hold multiple target formats like \texttt{jair} or \texttt{ieee\_transaction}.
Each template must contain a \texttt{paper.tex} file and a \texttt{Makefile}.
A background scanner loops through this directory and provides valid templates to the GUI.

The \texttt{PaperConverter} class copies the chosen template to \texttt{output/latex/} and injects data in three passes:

\begin{enumerate}[nosep]
	\item \textbf{Metadata:} String replacement swaps \texttt{\%\%TITLE\%\%} with the paper title. A parser scans for \texttt{\%\%BEGIN\_AUTHOR\%\%} and \texttt{\%\%END\_AUTHOR\%\%} markers. It duplicates the inner block for each author listed in the settings, replacing tags like \texttt{\{\{name\}\}} and \texttt{\{\{email\}\}} with user data.
	\item \textbf{Content:} The system maps each section to a template placeholder like \texttt{\%\%\allowbreak{}INTRODUCTION\allowbreak{}\%\%}. Rather than using rule-based text replacement, the \texttt{MarkdownToLaTeX.\allowbreak{}convert\_\allowbreak{}section\_\allowbreak{}to\_\allowbreak{}latex} method uses an LLM to translate the formatting into valid \LaTeX{} syntax. 
	A system prompt forces the model to format tables with booktabs, scale images to the text width, map citation keys to \texttt{\textbackslash cite\{citationKey\}} and shield special characters. 
	The \texttt{PaperConverter.\_inject\_sections\_into\_tex} method injects the formatted text strings.
	\item \textbf{Assets:} The \texttt{PaperConverter.\_copy\_plot\_images} method copies experiment plots from \texttt{output/experiments/plots/} directly into the target \texttt{images/} directory. This prevents broken \LaTeX{} figure references.
\end{enumerate}

Finally, the \texttt{PaperConverter.compile\_latex} method calls \texttt{make} inside the output directory.
This triggers the template's custom \texttt{Makefile} and builds the PDF.
A 60-second timeout aborts hanging \LaTeX{} compilations and drops an error flag.